% =============================================================================
% Introduction --- DNB variant (Datasets & Benchmarks track)
% Artifact-first framing; A/B/C claim grades; neutral wording.
% Target: <= 1 page.
% =============================================================================
\section{Introduction}
\label{sec:intro}

Reinforcement learning post-training now underpins much of modern
language-model reasoning and alignment~\citep{ouyang2022training,%
guo2025deepseekr1nature}, but empirical claims in this area are hard to
compare across stacks. PPO~\citep{schulman2017proximal},
GRPO~\citep{shao2024deepseekmath}, and DPO~\citep{rafailov2024direct} are
typically evaluated on different model families, with different framework
defaults, different reward functions, and different notions of success. The
field needs auditable, tier-classified artifacts more than it needs another
leaderboard.

\paragraph{Artifact-first contribution.}
\ours{} is released as a reproducible artifact: 79 runs spanning 7 RL
libraries, 5 model families, 32+ checkpoints, and total-parameter scales
from 0.6B to 671B (active-parameter counts 0.6B--37B for the largest
mixture-of-experts checkpoints), each accompanied by pinned containers,
step-level reward traces, ZVF/GU diagnostics, and (where the backend
permits) raw per-group rollout JSONL. Every (backend, model, task) cell
in the manifest is classified into one of three evidence tiers, used
consistently in the abstract, the main evidence table, and the
conclusion:
\begin{itemize}\setlength\itemsep{1pt}
  \item \textbf{(A-grade)} multi-seed completed cells used for inference;
        these support the only inferential claims in the paper.
  \item \textbf{(B-grade)} single-seed completed case-study runs on the
        closed-backend Tinker API, reported descriptively only.
  \item \textbf{(C-grade)} failed, partial, or never-started cells,
        enumerated in the run manifest for honesty.
\end{itemize}

\paragraph{What the artifact lets reviewers do.}
A reviewer can: (i) re-run any A-grade open-stack cell end-to-end via
\texttt{make reproduce-main}; (ii) inspect step-level ZVF/GU and reward
traces for every reported run, including B-grade Tinker case studies;
(iii) read off a per-run auditability matrix (Table~\ref{tab:artifact-card})
that lists, for each cell, whether raw rollouts, ZVF traces, checkpoints,
and held-out evaluation are released; and (iv) consult the claims-evidence
table (Table~\ref{tab:claims-tier}) to see exactly which artifact each
claim depends on and what its tier is.

\paragraph{Two narrow scientific take-aways.}
The artifact supports two A-grade results stated conservatively. First, the
mechanical identity $\mathrm{ZVF}(p, G) = p^{G} + (1{-}p)^{G}$ characterises
when binary rewards produce zero centered reward contrast in the
group-standardized runner; we use it as
a descriptive diagnostic, not as a standalone causal predictor of
downstream performance, because ZVF is mechanically coupled to reward
sparsity, group size, and baseline accuracy. Second, later, separate five-seed
Qwen3-8B GSM8K evidence yields only a $+1.3$pp delta over a fixed base;
$p{=}0.26$ is a one-sample seed-level test, not a paired-seed test or
reviewed-record evidence. Strong online reward curves on a closed backend do not by
themselves establish generalization gains. All B-grade single-seed
case-study runs across model scale are reported as illustrative.

\paragraph{Stratified ZVF diagnostics.}
A separate stratified diagnostic suite documents that the pooled
ZVF-vs-performance correlation is dominated by task identity: within
GSM8K alone the apparent monotone relationship does not survive
(Pearson $r=+0.40$, $p=0.09$ across the 19 Tinker GSM8K runs;
$r=+0.13$, $p=0.79$ across the 7 model-varying GSM8K cells). We
release these stratified analyses precisely so that follow-up work can
audit them against richer multi-seed corpora.
